%% kapa: exam=1 round=21 year=2010 period=1 form=A subject=civil_law subject_ko=민법 seq=21 number=21
%% source: Q-Net 감정평가사 기출문제 (cst00309) artlSeq=5024181 file=extracted/1008/5024181/2010년 제21회 감정평가사 1차시험 기출문제/1교시.hwp
%% conversion: kordoc 4.24.51 -> tools/gen1.py (auto); review=proofread-sonnet; flags=none
\begin{question}{21}
의사표시에 관한 설명으로 옳지 않은 것은? (다툼이 있으면 판례에 의함)
\begin{choices}
  \item 표의자는 의사표시를 발신한 후 그 도달 전에 의사표시를 철회할 수 있다.
  \item 채권양도의 통지는 관념의 통지이지만, 의사표시에 준하여 도달주의가 적용된다.
  \item 미성년자가 예외적으로 행위능력을 가지는 경우에는 그 한도에서 의사표시의 수령능력이 인정된다.
  \item 착오로 의사표시가 취소되면, 그 의사표시를 요소로 하는 재산상의 법률행위는 소급적으로 무효로 된다.
  \item 상대방이 의사표시의 내용을 확인하지 않은 채 정당한 사유없이 수령을 거절하였다면, 상대방이 그 내용을 알 수 없으므로 의사표시는 도달되지 않은 것으로 본다.
\end{choices}
\end{question}
\answer{5}
